\label{sec:hq-vae}


\section{Implementing HQ-VAE}
\label{sec:hq-1}

Liu et al. suggested implementing Hierarchically-Quansited Variational Autoencoder (HQ-VAE) \cite{HQVAE}.
The idea behind this is to split the communication into a high-level strategy and a low-level strategy, where the low-level strategy is conditioned on the high-level strategy.
By splitting the communication into macro-level and micro-level strategies, the model significantly reduces the computational burden associated with large, flat codebooks.

Common features are first extracted through a shared encoder:

\begin{equation}
    h_{\text{shared}} = Enc_{\theta_E}(x)
\end{equation}

The top level, responsible for macro-level strategies, projects these shared features to a vector:
\begin{equation}
    z_{\text{top}} = MLP_{\text{top}}(h_{\text{shared}})
\end{equation}

The bottom level, responsible for micro-level strategies, projects these shared features conditioned on the top level to a vector:
\begin{equation}
    z_{\text{bottom}} = MLP_{\text{bottom}}([h_{\text{shared}}, z_{\text{top}}])
\end{equation}

These vectors are quantised using one of our aforementioned methods (VQ-VAE, SQ-VAE).

\begin{equation}
    z_{q, top} = \text{Quantiser}(z_{top})
\end{equation}

\begin{equation}
    z_{q, bottom} = \text{Quantiser}(z_{bottom})
\end{equation}

Finally, the quantised vectors are concatenated to produce the final full communication vector $z_q$:

\begin{equation}
    z_q = [z_{q, top}, z_{q, bottom}]
\end{equation}


\section{HQ-VAE Hyperparameters}
\label{sec:hq-2}

The hyperparameters used for SQ-VAE and VQ-VAE are inappropriate for the AI Mother Tongue architecture with a discrete bottleneck using HQ-VAE.
For this, each codebook is given a vocabulary size of 16 (giving an effective vocabulary size of 256), allowing different levels of detail in different codebooks.
The code size also had to be increased from eight to sixteen.
In HQ-VAE, the code size is split across the codebooks (in a 1:2 ratio), so a total code size of eight would allocate only two dimensions to the top-level latent. 
With vMF SQ-VAE at the top level, a 2-dimensional unit sphere cannot support 16 codes, leading to a collapsed top-level codebook.
Using the code size of sixteen gives a sufficient five dimensions for the top-level codebook.


\section{HQ-VAE Analysis}
\label{sec:hq-3}

Replacing the single-level VAE bottleneck with an HQ-VAE did not significantly change task-level performance.
HQ-VAE's goal was for the top-level codebook to capture macro-level strategies and the bottom-level codebook to capture micro-level strategies.
To test this, I applied the same analyses as in \ref{sec:4.3} - \ref{sec:4.7} on each codebook and on the bottom-level codebook symbols conditioned on the top-level symbols.
In all cases, no clear, distinguishable strategies emerged at either level.
The lack of strategies in either codebook is likely due to the flat nature of the task (e.g. Find Station A).
A hierarchical language may require a more hierarchical task.